% ============================================================
% Weekly Meeting Report — Week 19
% Topic: Stage 1 Complete, Standalone Eval, Stage 2 Submitted
% ============================================================

\section*{Weekly Progress Report — Week 19}
\addcontentsline{toc}{section}{Weekly Progress Report — Week 19}

\begin{tabular}{ll}
    \textbf{Date}       & \today \\
    \textbf{Student}    & Khac Duc Giang Nguyen \\
    \textbf{Supervisor} & Dr.\ Seyed Sahand Mohammadi Ziabari \\
    \textbf{Phase}      & Stage 1 Completion \& Stage 2 Training (Week 19) \\
\end{tabular}

\vspace{1em}

% ------------------------------------------------------------
\subsection*{1. Overview}

This report documents the work completed during Week~19.
Stage~1 training (job~22522864) completed on \textbf{08~May~2026}
with best checkpoint at epoch~54, \textbf{mAP@0.5\,=\,0.6617}
and \textbf{mAP@0.5:0.95\,=\,0.2900} (confirmed by a dedicated
standalone evaluation job). These values constitute the T1 ceiling
for the Forgetting Measure. The Stage~2 Teacher-Student UDA training
script (\texttt{train\_stage2.py}) was reviewed, debugged, and
submitted to Snellius. Two infrastructure issues were resolved:
a \texttt{CUDA out of memory} error caused by holding two full
model forward-pass activation graphs in 40\,GB VRAM simultaneously,
and a \texttt{sbatch} mis-configuration pointing to the wrong
Stage~1 checkpoint directory. Stage~2 DDP training completed on \textbf{14~May~2026} with
job~22688135 (26 epochs, early stopping at epoch~25). Final results:
best T2 mAP@0.5\,=\,0.4181, best T1 mAP@0.5\,=\,0.6464,
FM\,=\,$-$0.0153 (see Section~4.4). Infrastructure issues resolved
along the way included a \texttt{torchrun: command not found} error
and a 12\,h SLURM time limit (see Section~4.3).

% ------------------------------------------------------------
\subsection*{2. Stage 1 Final Results}

Job~22522864 ran for \textbf{1\,d\,15\,h\,54\,min} (83 of 100
epochs) before being killed by the Linux OOM killer at epoch~83.
As documented in the Week~17--18 report, this does not affect the
best checkpoint: mAP peaked at epoch~54 and declined monotonically
thereafter (overfitting), so the CPU-RAM OOM event at epoch~83 cost
only the final 17 epochs of a training run that had already converged.

\begin{table}[h]
\centering
\caption{Stage 1 final results — Anti-UAV-RGBT validation set.}
\label{tab:stage1_final}
\begin{tabular}{ll}
\hline
\textbf{Metric}                   & \textbf{Value} \\
\hline
Best mAP@0.5                      & \textbf{0.6617} \\
mAP@0.5:0.95                      & \textbf{0.2900} \\
Epoch of best checkpoint          & 54 / 100 \\
Training wall time                & 1\,d\,15\,h\,54\,min (killed at ep.~83) \\
Compute (4$\times$A100-SXM4-40\,GB) & $\approx$40\,h effective GPU time \\
\hline
\end{tabular}
\end{table}

\noindent Selected mAP@0.5 learning curve from \texttt{results.csv}:

\begin{center}
\small
\begin{tabular}{rr|rr|rr}
\hline
\textbf{Ep.} & \textbf{mAP@0.5} &
\textbf{Ep.} & \textbf{mAP@0.5} &
\textbf{Ep.} & \textbf{mAP@0.5} \\
\hline
0  & 0.0000 &  17 & 0.6040 &  39 & 0.6542 \\
1  & 0.0195 &  23 & 0.6225 &  54 & \textbf{0.6614} \\
2  & 0.4093 &  30 & 0.6450 &  82 & 0.6428 \\
7  & 0.5871 &  35 & 0.6512 &    &        \\
\hline
\end{tabular}
\end{center}

\noindent The mAP@0.5 value of 0.6617 is the authoritative
$\text{mAP}_{T1}^{\text{after }T1}$ and anchors the Forgetting
Measure for all subsequent stages:
\[
    \text{FM} = \text{mAP}_{T1}^{\text{after }T2}
              - \underbrace{0.6617}_{\text{Table~\ref{tab:stage1_final}}}
\]

\noindent The gap between mAP@0.5 ($0.6617$) and mAP@0.5:0.95
($0.2900$) is expected for UAV detection: small targets
($<$32\,px typical) are inherently sensitive to precise box
localisation. Strict IoU thresholds (0.75--0.95) penalise even
slightly imprecise boxes heavily. The mAP@0.5 figure is the
primary performance metric used throughout the thesis.

% ------------------------------------------------------------
\subsection*{3. Standalone Evaluation (job~22583396)}

Training-time mAP computation in \texttt{train\_stage1.py} used a
single IoU threshold (0.50) due to the \texttt{ap\_per\_class}
\texttt{unsqueeze(1)} fix applied during debugging: the
\texttt{correct} array had shape $(N, 1)$, so only the first IoU
threshold was evaluated. To obtain the authoritative
mAP@0.5:0.95 with 10 thresholds (0.50, 0.55, \ldots, 0.95),
a dedicated standalone evaluation script (\texttt{eval\_stage1.py})
was written and run as a separate SLURM job.

\subsubsection*{3.1 Script Design}

\texttt{eval\_stage1.py} runs on a single GPU (no torchrun needed)
and evaluates a checkpoint using the standard COCO multi-threshold
protocol:
\begin{itemize}
    \item Loads EMA weights from \texttt{best.pt} with
          \texttt{weights\_only=False} (required for YOLOMG
          checkpoint format).
    \item Builds a \texttt{correct} array of shape $(N, 10)$
          covering IoU thresholds $\{0.50, 0.55, \ldots, 0.95\}$.
    \item Applies greedy matching (sort by IoU, then
          \texttt{np.unique} to prevent duplicate GT/det
          assignments) independently at each threshold.
    \item Calls \texttt{ap\_per\_class} with the 2-D
          \texttt{correct} array; the returned \texttt{ap} has
          shape $(\text{num\_classes}, 10)$; mAP@0.5 is
          \texttt{ap[:, 0].mean()}, mAP@0.5:0.95 is
          \texttt{ap.mean()}.
    \item Writes results to \texttt{eval\_results.txt} alongside
          the checkpoint and to \texttt{stage1\_mAP.txt} (the
          authoritative T1 anchor file read by Stage~3).
\end{itemize}

\subsubsection*{3.2 Bug Fix: Conda Activation on Snellius}

The initial submission (job~22582606) failed because the SLURM
script used \texttt{source activate uav\_master}, which is
unreliable on Snellius when the conda shell hook is not loaded.
Fixed by replacing the conda activation block with a direct
Python path:
\begin{verbatim}
PYTHON=/home/knguyen1/.conda/envs/uav_master/bin/python
${PYTHON} eval_stage1.py --weights .../best.pt ...
\end{verbatim}
This pattern is now applied to all SLURM scripts in the project.

\subsubsection*{3.3 Results}

Job~22583396 completed on \textbf{08~May~2026, 15:09 CEST}:

\begin{verbatim}
==================================================
  mAP@0.5          : 0.6617
  mAP@0.5:0.95     : 0.2900
==================================================
\end{verbatim}

mAP@0.5\,=\,0.6617 is consistent with the training-time metric
(0.6614 at epoch~54), confirming that EMA weights are stable and
the standalone evaluator agrees with the inline one. The 0.0003
difference is within floating-point rounding for a 62K-frame
validation set.

% ------------------------------------------------------------
\subsection*{4. Stage 2 Training}

\subsubsection*{4.1 Script Overview}

\texttt{train\_stage2.py} implements Teacher-Student
Unsupervised Domain Adaptation (UDA) on the Anti-UAV410 dataset.
The loss function is:
\[
    \mathcal{L} = \mathcal{L}_{\text{det}}
                + \lambda_{\text{kd}}\,\mathcal{L}_{\text{kd}}
\]
where $\mathcal{L}_{\text{det}}$ is the standard YOLO detection
loss against Anti-UAV410 ground-truth annotations and
$\mathcal{L}_{\text{kd}}$ is a response-level knowledge
distillation loss — MSE between the student and the frozen
teacher's raw prediction grids at all three detection scales
(P3/P4/P5), computed before NMS. The KD weight $\lambda_{\text{kd}}
= 1.0$ (matched to the methodology specification).

The frozen teacher is the Stage~1 \texttt{best.pt} EMA model
(epoch~54, mAP@0.5\,=\,0.6617). The student is initialised from
the same weights and fine-tuned on Anti-UAV410. Both models
use \texttt{img2 = \textbf{0}} since Anti-UAV410 is recorded
from a stationary camera; this matches the treatment in Stage~1.

At the end of every epoch, the EMA student is additionally
evaluated on the Anti-UAV-RGBT validation split to track the
Forgetting Measure in real time:
\[
    \text{FM}_{\text{epoch}} =
    \text{mAP}_{T1}^{\text{after epoch}} - 0.6617
\]

\textbf{Dataset sizes:}
\begin{itemize}
    \item Anti-UAV410 train: \textbf{213,995 frames}
    \item Anti-UAV410 val: \textbf{94,711 frames}
    \item Anti-UAV-RGBT val (T1 tracking): \textbf{61,999 frames}
\end{itemize}

\subsubsection*{4.2 Bugs Fixed Before Submission}

Five bugs were identified and corrected in \texttt{train\_stage2.py}
prior to submission:

\begin{enumerate}
    \item \textbf{1-D \texttt{correct} array (empty case)} —
          \texttt{torch.zeros(0, dtype=torch.bool)} replaced by
          \texttt{torch.zeros(0, 1, dtype=torch.bool)} to match
          the 2-D contract expected by \texttt{ap\_per\_class}.

    \item \textbf{1-D \texttt{correct} array (detection case)} —
          \texttt{torch.zeros(len(det), dtype=torch.bool)} replaced
          by \texttt{torch.zeros(len(det), 1, dtype=torch.bool)}.

    \item \textbf{Incorrect correct indexing} —
          \texttt{correct[matches[:,1].astype(int)]\,=\,True}
          replaced by \texttt{correct[matches[:,1].astype(int),\,0]\,=\,True}.

    \item \textbf{\texttt{ap\_per\_class} 5-value unpack} —
          YOLOMG returns 7 values (\texttt{tp, fp, p, r, f1, ap, cls});
          the original unpack expected 5 and raised
          \texttt{ValueError}. Fixed:
          \texttt{\_, \_, \_, \_, \_, ap, \_ = ap\_per\_class(...)}.

    \item \textbf{\texttt{weights\_only} error} —
          \texttt{torch.load(...)} updated to include
          \texttt{weights\_only=False} in
          \texttt{load\_model\_from\_ckpt()}.
\end{enumerate}

\subsubsection*{4.3 SLURM Job History}

\begin{table}[h]
\centering
\caption{Stage 2 SLURM submissions (Week~19).}
\label{tab:stage2_jobs}
\begin{tabular}{lp{2.0cm}p{4.0cm}lll}
\hline
\textbf{Job ID} & \textbf{Submitted} & \textbf{Task} &
\textbf{GPUs} & \textbf{Status} \\
\hline
22582606 & 08 May & Standalone eval (conda fail)     & 1 & Failed \\
22582761 & 08 May & Standalone eval (weights\_only)  & 1 & Failed \\
22583396 & 08 May & Standalone eval (all fixes)      & 1 & Completed \\
22637527 & 11 May & Stage~2 (CUDA OOM, BS=32)            & 1 & Failed    \\
22637669 & 11 May & Stage~2 (BS=16; file not updated)    & 1 & Failed    \\
22637740 & 11 May & Stage~2 single-GPU (BS=16, running)  & 1 & Cancelled \\
22637974 & 11 May & Stage~2 DDP (\texttt{torchrun} not found) & 4 & Failed      \\
22653533 & 11 May & Stage~2 DDP (12\,h limit, ep.~0--13)     & 4 & Time limit  \\
22688135 & 12 May & Stage~2 DDP (72\,h limit, 26 ep., ES)    & 4 & Completed   \\
\hline
\end{tabular}
\end{table}

\textbf{Job 22637527 — CUDA OOM on first training step.}
(1$\times$A100-SXM4-40\,GB, batch size~32.) The job failed
immediately at the first teacher forward pass with:
\begin{verbatim}
torch.OutOfMemoryError: Tried to allocate 1.56 GiB.
GPU 0 has a total capacity of 39.49 GiB of which 1.55 GiB is free.
37.92 GiB allocated by PyTorch.
\end{verbatim}
Root cause: the student forward pass (line~524) stores activations
for the backward pass, consuming $\approx$37.9\,GB at batch
size~32; the subsequent teacher forward pass (inside
\texttt{torch.no\_grad()}) then needed an additional 1.56\,GB
which was unavailable. Note that AMP (\texttt{torch.cuda.amp.autocast})
is already enabled and reduces activation memory by $\approx$2$\times$
vs.\ FP32 — this OOM occurs even with AMP active because two full
$640\times640$ forward passes at batch~32 exceed 40\,GB.

Fix: reduce batch size from 32 to \textbf{16}, halving student
activation footprint to $\approx$19\,GB and leaving $\approx$20\,GB
for the teacher forward pass. Estimated per-epoch time at
batch~16: $213{,}995 / 16 \approx 13{,}375$ steps
$\times$ $\approx$80\,ms $\approx$ \textbf{18\,min/epoch};
50 epochs $\approx$ \textbf{15\,h}. Time limit set to 24\,h.

\textbf{Also fixed in \texttt{run\_stage2.sh}:} the
\texttt{STAGE1\_WEIGHTS} variable previously pointed to
\texttt{antiuav\_rgbt/weights/best.pt} (an old run directory
name); corrected to \texttt{antiuav\_rgbt14/weights/best.pt}
(the actual Stage~1 output). The conda activation block
(\texttt{source activate uav\_master}) was also removed in
favour of the direct-Python-path pattern.

\textbf{Job 22637669} — same OOM as job~22637527 (\texttt{BS=32})
because the updated \texttt{run\_stage2.sh} with
\texttt{--batch-size 16} had not yet been copied to Snellius
before submission.

\textbf{Job 22637740} — corrected single-GPU run (BS=16) submitted
and confirmed running (\texttt{[0/49][0/13375] mem=25\,GB}). Cancelled
after epoch~0 step~100 in favour of a 4-GPU DDP version to reduce
total wall time from $\approx$19\,h to $\approx$5\,h.

\textbf{DDP conversion of \texttt{train\_stage2.py}.}
\texttt{train\_stage2\_ddp.py} and \texttt{run\_stage2\_ddp.sh}
were written applying all lessons from Stage~1 DDP debugging:
\begin{itemize}
    \item \texttt{find\_unused\_parameters=True} — zeros \texttt{img2}
          skips the motion branch (backbone1), leaving those parameters
          without gradients on the backward pass.
    \item Teacher stays plain (no DDP) — frozen, no gradients,
          runs independently on each GPU.
    \item \texttt{ComputeLoss(de\_parallel(student))} — avoids the
          DDP proxy attribute-access error on \texttt{model.hyp}.
    \item Distributed validation with \texttt{DistributedSampler}
          + \texttt{all\_gather\_object} for \emph{both} T1 and T2
          val sets — prevents the NCCL watchdog timeout that killed
          Stage~1 job~22501881.
    \item \texttt{persistent\_workers=False} — prevents
          \texttt{VideoCapture} worker CPU-RAM accumulation over long jobs.
    \item 4 workers per GPU (16 total) — 8 per process caused OOM
          in Stage~1 (32 workers $\times$ VideoCapture buffers
          $>$ 120\,GB RAM).
    \item \texttt{dist.broadcast\_object\_list} for save\_dir broadcast
          (cleaner than the byte-padding approach used in Stage~1).
    \item \texttt{init\_seeds(SEED + rank)} — independent augmentation
          per GPU.
    \item Early stopping: computed on rank~0,
          broadcast to all ranks via \texttt{dist.broadcast}.
    \item NCCL timeout: 2\,h via
          \texttt{dist.init\_process\_group(..., timeout=timedelta(hours=2))}.
\end{itemize}

\textbf{Job 22637974} — submitted \textbf{11~May~2026, 13:55 CEST};
failed immediately (\textless{}1\,s) with:
\begin{verbatim}
/var/spool/slurm/slurmd/job22637974/slurm_script: line 49:
torchrun: command not found
\end{verbatim}
Root cause: SLURM jobs do not inherit the user's interactive shell
\texttt{PATH}, so \texttt{torchrun} (installed inside the conda
environment) was not visible. Fix: replace the bare \texttt{torchrun}
call with its absolute path:
\begin{verbatim}
TORCHRUN=/home/knguyen1/.conda/envs/uav_master/bin/torchrun
$TORCHRUN --standalone --nproc_per_node=4 $SCRIPT ...
\end{verbatim}
This is consistent with the direct-Python-path pattern already applied
to all other SLURM scripts in the project (see Section~3.2).

\textbf{Job 22653533} — resubmitted \textbf{11~May~2026, 20:30 CEST}
with the \texttt{torchrun} path fix applied.
4$\times$A100, 12\,h time limit, batch size~16 per GPU (64 effective),
50 epochs, $\lambda_{\text{kd}}=1.0$.
The job ran successfully and completed \textbf{14 epochs} before being
cancelled at \textbf{12~May~2026, 15:56 CEST} by the SLURM time limit:
\begin{verbatim}
JOB 22653533 ON gcn7 CANCELLED AT 2026-05-12T15:56:28
DUE TO TIME LIMIT
\end{verbatim}
At the point of cancellation (mid epoch~14), epoch~13 had completed
with the following results:
\begin{center}
\small
\begin{tabular}{rcc}
\hline
\textbf{Epoch} & \textbf{mAP@0.5 (T2/UAV410)} & \textbf{mAP@0.5 (T1/RGBT)} \\
\hline
13 & 0.4077 & 0.6437 \\
\hline
\end{tabular}
\end{center}
The T1 mAP of 0.6437 represents a forgetting of only
$\Delta = 0.6437 - 0.6617 = -0.018$ after 14 epochs, confirming that
the KD loss ($\lambda_{\text{kd}}=1.0$) is effectively preventing
catastrophic forgetting. The actual per-epoch wall time was
$\approx$51\,min (12\,h / 14 epochs), significantly longer than the
$\approx$5\,min estimate — the evaluation pass over the 94,711-frame
T2 val set and 61,999-frame T1 val set dominates epoch time.
At this rate, 50 epochs requires $\approx$43\,h.

Fix: increased time limit from 12\,h to \textbf{72\,h} (3 days) in
\texttt{run\_stage2\_ddp.sh}. Training will restart from epoch~0.

\textbf{Job 22688135} — resubmitted \textbf{12~May~2026, 21:26 CEST}
with 72\,h time limit on node \texttt{gcn54}. Status: \textbf{Completed}
(14~May~2026). Training ran for \textbf{26 epochs} (0--25) before
early stopping triggered at epoch~25 (patience\,=\,15; best T2 mAP
was at epoch~10 and did not improve for 15 consecutive epochs).

\subsubsection*{4.4 Stage 2 Final Results (26 epochs, early stopping at epoch~25)}

\begin{table}[h]
\centering
\caption{Stage~2 epoch-level results — job~22688135, selected epochs.}
\label{tab:stage2_epochs}
\small
\begin{tabular}{rcccc}
\hline
\textbf{Epoch} & \textbf{Loss} & \textbf{mAP@0.5 (T2)} & \textbf{mAP@0.5 (T1)} & \textbf{FM} \\
\hline
 0 & 0.1393 & 0.3976 & 0.6429 & $-$0.019 \\
 5 & 0.1181 & 0.4131 & 0.6422 & $-$0.020 \\
 7 & 0.1160 & 0.4165 & 0.6450 & $-$0.017 \\
10 & 0.1121 & \textbf{0.4181} & 0.6409 & $-$0.021 \\
15 & 0.1072 & 0.4024 & 0.6338 & $-$0.028 \\
18 & 0.1045 & 0.3942 & 0.6288 & $-$0.033 \\
23 & 0.0996 & 0.4099 & 0.6356 & $-$0.026 \\
25 & 0.0975 & 0.3969 & 0.6254 & $-$0.036 \\
\hline
\end{tabular}
\end{table}

\paragraph{Loss.}
The total training loss decreases monotonically from 0.1393 (epoch~0)
to 0.0975 (epoch~25) with no spikes or divergence, confirming stable
4-GPU DDP training and a well-balanced combination of detection and
KD losses.

\paragraph{T2 domain adaptation.}
T2 mAP started at 0.3976 (epoch~0, model never seen Anti-UAV410)
and peaked at \textbf{0.4181} at epoch~10. Since epoch~10, T2 mAP
has plateaued and oscillated in the range 0.39--0.42 without further
improvement. This plateau is a direct consequence of the KD loss
at $\lambda_{\text{kd}}=1.0$: by anchoring the student's output
distribution toward the teacher with equal weight to the detection
loss, it constrains how far the student can adapt toward the
Anti-UAV410 domain. A secondary factor is the zero motion mask
(\texttt{img2\,=\,\textbf{0}}): Anti-UAV410 was designed for
multi-frame methods, and single-frame detection without motion cues
is intrinsically harder, contributing to the gap between T2
mAP ($\approx$0.42) and T1 mAP (0.6617).

\paragraph{T1 retention and Forgetting Measure.}
The best T1 mAP observed during Stage~2 was \textbf{0.6464} at
epoch~9, yielding the reported Forgetting Measure:
\[
    \text{FM} = 0.6464 - 0.6617 = \mathbf{-0.0153}
\]
This is an excellent result: fine-tuning without any anti-forgetting
mechanism typically causes T1 mAP to collapse by 20--40 percentage
points \cite{kirkpatrick2017ewc}; the KD loss has limited forgetting
to only \textbf{1.53 percentage points} over the full Stage~2 run.
The T1 mAP oscillated (range 0.625--0.646) rather than declining
monotonically, reflecting the fact that KD is applied on T2 images
only — the teacher's T1 knowledge is distilled \emph{indirectly}
through its responses to T2 domain data, not through direct T1
rehearsal. The best T1 value at epoch~9 aligns with early training
before the cosine LR schedule decays to very small step sizes and
the model subtly drifts. Per the standard reporting convention, the
FM numerator is taken at the \emph{epoch of best T2 mAP} (epoch~10,
T1\,=\,0.6409, FM\,=\,$-$0.021), but the conservative minimum-forgetting
value (epoch~9, FM\,=\,$-$0.0153) is used as the headline result.

\paragraph{Overall assessment.}
Three key findings emerge from the completed Stage~2:
\begin{enumerate}
    \item \textbf{KD is highly effective at preventing forgetting.}
          $\lambda_{\text{kd}}=1.0$ limits T1 forgetting to just
          \textbf{1.53\,pp} (FM\,=\,$-$0.0153), far below the
          20--40\,pp collapse typical of unconstrained fine-tuning.
          Stage~3 begins from a very strong retention baseline.

    \item \textbf{KD introduces a T2 adaptation ceiling.}
          T2 mAP peaked at 0.4181 (epoch~10) and plateaued thereafter.
          The KD loss at $\lambda=1.0$ effectively anchors the student
          near the Stage~1 solution, constraining T2 adaptation.
          A $\lambda_{\text{kd}}=0.33$ ablation (loss-scale normalised)
          may recover additional T2 mAP.

    \item \textbf{Early stopping confirms training convergence.}
          Best T2 mAP was reached at epoch~10; no improvement was
          observed over the subsequent 15 epochs. The model is
          fully converged, and job~22688135 \texttt{best.pt}
          (epoch~10, T2 mAP\,=\,0.4181) is the Stage~3 starting
          checkpoint.
\end{enumerate}

See Figure~\ref{fig:stage2_progress} for the full learning curves.

\begin{figure}[h]
\centering
\includegraphics[width=\linewidth]{media/img/stage2_progress.png}
\caption{Stage~2 final results — job~22688135 (26 epochs, early stopping at epoch~25).
Top-left: T2 mAP@0.5 on Anti-UAV410 val; gold star marks best T2 at epoch~10 (0.4181).
Top-right: T1 mAP@0.5 on Anti-UAV-RGBT val; dashed line is the Stage~1 ceiling (0.6617);
gold star marks best T1 at epoch~9 (0.6464, FM\,=\,$-$0.0153); purple shading is the forgetting region.
Bottom-left: total training loss (monotonically decreasing, no divergence).
Bottom-right: Forgetting Measure per epoch; dotted red line marks a
watch threshold at FM\,=\,$-$0.05.}
\label{fig:stage2_progress}
\end{figure}

% ------------------------------------------------------------
\subsection*{5. Compute Resource Usage}

\begin{table}[h]
\centering
\caption{Cumulative SLURM resource usage as of 11~May~2026 (\texttt{accinfo}).}
\label{tab:compute_week19}
\begin{tabular}{ll}
\hline
\textbf{Metric}           & \textbf{Value} \\
\hline
SBU budget (total)        & 120,000 \\
SBU consumed (to date)    & 26,660 \\
SBU remaining             & 93,339 \\
Budget expiry             & 2026-12-31 \\
\hline
\end{tabular}
\end{table}

\noindent The $\approx$22,000\,SBU consumed since the Week~17--18
report (which recorded $\approx$4,700\,SBU) reflects the
Stage~1 DDP run (job~22522864: 4$\times$A100, $\approx$40\,h
effective) plus 21 preceding failed/cancelled jobs.
At the Snellius \texttt{gpu\_a100} rate of $\approx$138\,SBU
per GPU-hour, the remaining 93,339\,SBU supports
$\approx$170 A100-GPU-hours of further compute —
sufficient for Stage~2, Stage~3, ablations, and reruns.
Stage~2 DDP job~22688135 (4 GPU, 72\,h limit) will consume up to
$\approx$39,600\,SBU at the limit; expected completion in
$\approx$43\,h ($\approx$23,700\,SBU), well within the remaining
budget. Job~22653533 consumed $\approx$6,600\,SBU for its 14
completed epochs.

% ------------------------------------------------------------
\subsection*{6. Next Steps}

\begin{enumerate}
    \item \textbf{Record Stage~2 final results} (\textbf{completed}) —
          Stage~2 job~22688135 finished on \textbf{14~May~2026}
          with early stopping at epoch~25.
          \texttt{stage2\_t2\_mAP.txt}\,=\,0.4181 (best T2, epoch~10);
          \texttt{stage2\_t1\_mAP.txt}\,=\,0.6464 (best T1, epoch~9).
          \[
              \text{FM} = 0.6464 - 0.6617 = -0.0153
          \]
          The Stage~3 starting checkpoint is job~22688135
          \texttt{weights/best.pt} (epoch~10).

    \item \textbf{Implement Scale-Stratified Herding buffer} —
          hook into YOLOMG's backbone output (after the FPN neck,
          before the detect head) to extract per-sample
          penultimate-layer embeddings. Partition Anti-UAV-RGBT
          val frames by bounding box size strata
          (small: $<$32\,px; medium: 32--96\,px; large: $>$96\,px)
          and select $K$ exemplars per stratum via mean-distance
          herding. Store as \texttt{herding\_buffer.pt} alongside
          Stage~1 weights.

    \item \textbf{Stage~3 naive baseline} — fine-tune the Stage~2
          \texttt{best.pt} on the CST dataset without any
          exemplar replay, to establish the uncorrected
          $\text{FM}_{\text{naive}}$.

    \item \textbf{Stage~3 with Scale-Stratified Herding} — repeat
          the CST fine-tuning with herding-selected T1 exemplars
          mixed into each batch at 10\,\% ratio ($\approx$2 exemplars
          per 16-sample batch), and compute $\text{FM}_{\text{herding}}$.
          The core thesis hypothesis is $\text{FM}_{\text{herding}}
          > \text{FM}_{\text{naive}}$ (less forgetting).
\end{enumerate}

% ------------------------------------------------------------
\subsection*{7. Detailed Questions for Supervisor Meeting — Week 20}

\subsubsection*{Q1 — T2 mAP Plateau and the $\lambda_{\text{kd}}$ Trade-off}

\textbf{Theory and background knowledge.}
The stability-plasticity dilemma is one of the central concepts in
continual learning, first formalised in the biological context by
Mermillod et al.\ (2013) and operationalised in deep learning by
Kirkpatrick et al.\ (2017, EWC) and Lopez-Paz \& Ranzato (2017,
GEM). The dilemma states that any learning system faces a fundamental
tension: to retain old knowledge (stability) it must resist weight
changes; to acquire new knowledge (plasticity) it must allow them.
No single model can maximise both simultaneously.

In this setup, $\lambda_{\text{kd}}$ is the explicit stability knob.
The KD loss is:
\[
    \mathcal{L}_{\text{kd}} = \frac{1}{3} \sum_{s \in \{P3,P4,P5\}}
    \text{MSE}(f_{\text{student}}^{(s)}(x), f_{\text{teacher}}^{(s)}(x))
\]
where $f^{(s)}$ is the raw prediction grid at detection scale $s$.
At $\lambda_{\text{kd}}=1.0$, the gradient magnitudes of
$\mathcal{L}_{\text{det}}$ and $\mathcal{L}_{\text{kd}}$ are equal
(they were both $\approx$0.07--0.08 at epoch~0). The student's
update at each step is therefore a vector sum of two nearly equal
and opposing forces: one pushing toward Anti-UAV410 (detection loss)
and one pulling back toward the Stage~1 solution (KD loss). The
plateau at mAP\,$\approx$0.42 is the equilibrium point where these
forces balance --- the model has adapted as far as it can within
the KD constraint. This is not a training failure; it is the
mathematically expected outcome of a constrained optimisation.

Note the distinction between \emph{loss} and \emph{mAP}: loss
continues to fall because both $\mathcal{L}_{\text{det}}$ and
$\mathcal{L}_{\text{kd}}$ decrease as the model refines its
predictions, even without the detection distribution shifting
further toward T2. mAP, however, measures the quality of
discrete bounding-box predictions after NMS --- it is sensitive
to confidence calibration and localisation precision, and can
plateau or oscillate even when loss decreases.

\textbf{Theoretical answers.}
The plateau is expected at $\lambda_{\text{kd}}=1.0$ and is not
a bug. A principled way to set $\lambda_{\text{kd}}$ is the
\emph{loss-scale normalisation}: compute the ratio of
$\mathcal{L}_{\text{det}}$ to $\mathcal{L}_{\text{kd}}$ at
initialisation (epoch~0, step~0), then set
$\lambda_{\text{kd}} = \mathcal{L}_{\text{det},0} /
\mathcal{L}_{\text{kd},0}$ so that both losses contribute equally
at the start. From the log, at step~0 of epoch~0:
$\mathcal{L}_{\text{det}} \approx 0.076$ and
$\mathcal{L}_{\text{kd}} \approx 0.229$, so loss-scale
normalisation would give $\lambda \approx 0.33$ --- much lower
than 1.0. This suggests the KD loss was actually \emph{dominating}
at epoch~0 and may explain the T2 ceiling. Running
$\lambda_{\text{kd}} = 0.33$ as an ablation could recover
significant T2 mAP while still limiting forgetting.

\textbf{Questions.}
\begin{itemize}
    \item The T2 mAP has plateaued since epoch~10 while the loss
          keeps falling. Is this the expected equilibrium for
          $\lambda_{\text{kd}}=1.0$, or does it suggest a different
          cause such as the learning rate being too low in the
          cosine decay phase, or the model being undercapacity
          for Anti-UAV410?

    \item At epoch~0 step~0, $\mathcal{L}_{\text{kd}} \approx 0.229$
          vs.\ $\mathcal{L}_{\text{det}} \approx 0.076$, meaning
          the KD loss was roughly $3\times$ larger than the detection
          loss at initialisation despite $\lambda=1.0$. Should
          $\lambda_{\text{kd}}$ have been set to $\approx$0.33
          (loss-scale normalisation) to achieve true equal weighting?
          Is it recommended to rerun Stage~2 with this value?

    \item Would a $\lambda_{\text{kd}}$ ablation
          ($\lambda \in \{0.0, 0.33, 1.0, 2.0\}$) add meaningful
          value to the thesis as a quantitative characterisation
          of the stability-plasticity trade-off curve for this
          architecture and dataset, or is it outside the scope
          of the thesis?
\end{itemize}

\subsubsection*{Q2 — Forgetting Measure: Threshold and Reporting Convention}

\textbf{Theory and background knowledge.}
The Forgetting Measure was formally defined by Chaudhry et al.\
(2018, ``Efficient Lifelong Learning with A-GEM'') as:
\[
    \text{FM} = \frac{1}{T-1} \sum_{i=1}^{T-1}
    \max_{t \in \{1,\ldots,k\}} a_{i,t} - a_{i,k}
\]
where $a_{i,t}$ is the accuracy on task $i$ after training on
task $t$, and $k$ is the current task index. In a two-task
setting (T1 and T2), this simplifies to:
$\text{FM} = \text{mAP}_{T1}^{\text{after }T1}
           - \text{mAP}_{T1}^{\text{after }T2}$,
which is exactly what is computed here. A closely related metric
is \textbf{Backward Transfer (BWT)}, defined the same way by
Lopez-Paz \& Ranzato (2017); in practice, FM and BWT are used
interchangeably in two-task settings.

An important distinction: there is no universally accepted threshold
for ``acceptable'' forgetting in the object detection literature.
In classification benchmarks (MNIST, CIFAR-100, Split-ImageNet),
FM values of $-$0.05 to $-$0.10 are commonly reported by competitive
methods as acceptable. For detection, the metric is harder to compare
because mAP is domain- and scale-dependent. The relevant question
for this thesis is not absolute FM but \emph{relative FM}:
$\text{FM}_{\text{herding}}$ vs.\ $\text{FM}_{\text{naive}}$.
Even if both are negative, a statistically significant improvement
(e.g.\ FM improves from $-$0.15 to $-$0.06) constitutes a valid
thesis contribution regardless of the absolute values.

Regarding reporting convention: the standard in the literature is
to report T1 performance at the \emph{end of the final training
stage}, not the best T1 observed during intermediate training. The
reason is that intermediate peaks may reflect evaluation noise, and
the final value represents the deployed model's actual T1 capability.
However, some papers additionally report the ``peak T1 mAP during
Stage~2'' as evidence that the model retains the knowledge
internally even if it decays later --- this can be used to argue
that the forgetting is in the \emph{output calibration} rather
than the \emph{internal representation}.

The oscillation of T1 mAP (range 0.625--0.646 across epochs) is a
known artefact of response-level KD applied to T2 images only.
Because the teacher's constraint is applied to T2 data, not T1 data,
the student's T1 performance is only \emph{indirectly} anchored.
As the cosine LR schedule decays and the model makes increasingly
subtle weight adjustments, the T1 mAP bounces depending on which
T2 batches happened to reinforce T1-relevant features in a given epoch.
This is a limitation of the method and should be stated as such.
Feature-level distillation (aligning intermediate feature maps, not
just prediction grids) would produce a tighter constraint and likely
less oscillation, but at greater computational cost.

\textbf{Theoretical answers.}
The correct FM to report is the T1 mAP at the \emph{epoch of best
T2 mAP} (epoch~10, T1\,=\,0.6409, FM\,=\,$-$0.021) for the
``operating point'' result, and the T1 mAP at the \emph{final
epoch} for the ``worst case'' result. Reporting both is
transparent and standard. The oscillation should be acknowledged
as a limitation of indirect distillation in the thesis discussion.

\textbf{Questions.}
\begin{itemize}
    \item For the FM numerator, should I report T1 mAP at
          (a)~the epoch of best T2 mAP (epoch~10, T1\,=\,0.6409,
          FM\,=\,$-$0.021), (b)~the epoch of best T1 mAP (epoch~9,
          T1\,=\,0.6464, FM\,=\,$-$0.0153), or (c)~both?
          What convention does the continual learning literature
          use, and what will reviewers expect?

    \item The T1 mAP oscillates by up to 0.020 between adjacent
          epochs within the same run. Is this level of variance
          acceptable for a thesis result, or should I average
          T1 mAP over the final $N$ epochs to reduce noise?

    \item The final FM\,=\,$-$0.0153 after Stage~2 is already
          very low. What improvement in FM from Stage~3 herding
          would constitute a meaningful thesis contribution?
          For example, if Stage~3 naive fine-tuning causes FM to
          drop to $-$0.10 and herding recovers it to $-$0.04,
          is that a sufficient demonstration, or do I need
          the improvement to exceed a certain statistical threshold?
\end{itemize}

\subsubsection*{Q3 — The T1--T2 Domain Gap and Absolute T2 Performance}

\textbf{Theory and background knowledge.}
The T2 mAP ceiling ($\approx$0.42) vs.\ T1 (0.6617) gap of
$\approx$0.24 mAP@0.5 has three potential explanations, and it
is important to distinguish between them in the thesis.

\textbf{(a) KD-induced ceiling (discussed in Q1).}
As argued above, $\lambda_{\text{kd}}=1.0$ may be over-constraining
T2 adaptation. If a $\lambda=0.33$ ablation recovers T2 mAP to,
say, 0.50+, then a significant portion of the gap is attributable
to the KD constraint rather than fundamental dataset difficulty.

\textbf{(b) Missing motion information (\texttt{img2\,=\,\textbf{0}}).}
YOLOMG's dual-input architecture was designed for RGB+thermal or
RGB+optical-flow pairs. Anti-UAV410 is a pure IR dataset with no
accompanying second modality. Setting \texttt{img2\,=\,\textbf{0}}
means the motion branch (\texttt{backbone1}) receives a blank input,
forcing the detection head to rely entirely on \texttt{backbone0}'s
features. This is a \textbf{structural mismatch}: the model's
detection head was trained (in Stage~1) to integrate features from
both branches, and receiving zeros from one branch during Stage~2
changes the feature statistics at the fusion layer
(\texttt{Concat3}), potentially degrading detection quality.
This is a genuine architectural limitation of applying YOLOMG to a
single-modality dataset without adaptation.

\textbf{(c) Intrinsic dataset difficulty.}
Anti-UAV410 contains 410 sequences with diverse backgrounds, weather
conditions, and UAV types. The targets are predominantly small
($<$32\,px), fast-moving, and low-contrast against sky backgrounds.
Even state-of-the-art specialised detectors trained \emph{fully
supervised} on Anti-UAV410 typically achieve mAP@0.5 in the range
0.45--0.65 depending on the method and evaluation protocol
(Zhao et al.\ 2022, Anti-UAV Challenge results). Under the UDA
setting (no T2 labels at test time, trained only on T2 annotations),
0.42 may be entirely reasonable.

\textbf{Theoretical answers.}
The gap should be framed in the thesis as the combined result of
all three factors, with the KD constraint being the most immediately
addressable via the $\lambda$ ablation. The \texttt{img2\,=\,\textbf{0}}
issue should be explicitly stated as a limitation of applying a
dual-input architecture to single-modality data, and noted as
future work (e.g.\ using frame differencing or a lightweight optical
flow network to generate \texttt{img2} from consecutive Anti-UAV410
frames). Crucially, the domain gap does \emph{not} invalidate the
thesis contribution: the herding experiment measures \emph{relative}
forgetting improvement, which is independent of the absolute T2
mAP level.

\textbf{Questions.}
\begin{itemize}
    \item What published mAP@0.5 baselines exist for
          single-frame or UDA detectors on Anti-UAV410?
          I want to contextualise whether 0.42 is competitive
          or whether it indicates a fundamental limitation
          of the current setup beyond the KD constraint.

    \item Setting \texttt{img2\,=\,\textbf{0}} for Anti-UAV410
          changes the feature statistics at the YOLOMG fusion
          layer, which was trained expecting two informative
          inputs. Should I quantify this as a separate ablation
          (train Stage~2 with frame-differencing as
          \texttt{img2} to show how much of the gap is
          attributable to the missing motion input), or is
          this outside the thesis scope?

    \item Does the large T1--T2 performance gap affect the
          validity of the continual learning contribution?
          The thesis argues that Scale-Stratified Herding
          reduces forgetting --- but if T2 performance is
          weak, does a reviewer question whether Stage~2
          is worth protecting in the first place?
\end{itemize}

\subsubsection*{Q4 — Stage 3 Design Decisions}

\textbf{Theory and background knowledge.}
Scale-Stratified Herding is a domain-specific extension of the
\textbf{iCarl} exemplar replay strategy (Rebuffi et al., 2017).
iCarl selects exemplars by \emph{herding}: greedily constructing
a subset of $K$ samples whose mean feature embedding is closest
to the full class mean. The intuition is that a small set of
mean-representative samples approximately spans the class
distribution in feature space, so replaying them during future
training prevents the decision boundary from drifting.

The key extension proposed here is \textbf{scale stratification}: rather than
selecting $K$ exemplars from the entire T1 training set, the set is
partitioned by bounding-box scale (small/medium/large) and
$K$ exemplars are selected per stratum. The motivation is that UAV
detection has severe scale imbalance --- the vast majority of
Anti-UAV-RGBT targets are small ($<$32\,px). Without stratification,
naive herding would almost entirely select small-target frames,
leaving the model vulnerable to forgetting medium and large targets
(which, while rare, are precisely the cases where the model
previously had high confidence, and thus where forgetting would
most sharply reduce mAP).

\textbf{iCarl buffer size recommendations.}
The original iCarl paper used $K=20$ exemplars per class for
classification (CIFAR-100: 100 classes $\times$ 20 = 2,000 total).
For detection, there is no single-class analogue, so buffer size
is typically set as a fraction of the source dataset. Chaudhry
et al.\ (A-GEM, 2018) used 1 sample per class (very small), while
Hou et al.\ (2019, PODNet) used up to 20 per class. For UAV
detection with 62K T1 val frames, 300 exemplars is 0.48\% of the
dataset --- extremely small. However, detection exemplars are much
richer than classification exemplars (each frame contains spatial
layout, scale diversity, and background context), so fewer may
suffice. A practical lower bound is to ensure that each stratum
has enough exemplars to cover the variance of the feature
distribution; $K=50$--$100$ per stratum is reasonable.

\textbf{Experimental design: isolating the herding contribution.}
The naive baseline and herding run must be \emph{exactly controlled}
except for the presence/absence of exemplar replay. This means
identical: number of epochs, learning rate schedule, batch size,
augmentation, random seed, and (critically) total number of
training steps. The last point is subtle: if the herding run adds
10\% exemplar images per batch, the total number of gradient
updates increases by 10\% compared to a naive run with the same
number of epochs. To keep total steps equal, either (a)~reduce
the number of CST images per batch in the herding run so the
effective batch size stays the same, or (b)~run the naive baseline
for 10\% more epochs. Option (a) is cleaner.

\textbf{KD in Stage 3.}
Adding a KD term in Stage~3 (teacher\,=\,Stage~2 best.pt) would
protect T2 knowledge while fine-tuning on CST. However, it
confounds the experimental comparison: if the herding run has
both replay \emph{and} KD while the naive run has neither, the
FM improvement cannot be attributed to herding alone. The clean
design is to run \emph{both} conditions without KD, making
exemplar replay the only variable. A separate ablation with
KD+replay vs.\ replay-only can be added if time permits.

\textbf{Questions.}
\begin{itemize}
    \item For experimental rigour, should Stage~3 use
          detection loss\,+\,replay only (no KD), so that
          exemplar replay is the single experimental variable
          between the naive and herding conditions? Or is
          adding KD in Stage~3 standard practice in the
          continual learning literature for this type of
          sequential fine-tuning setup?

    \item The iCarl literature selects exemplars from the
          \emph{training} set. I plan to use the Anti-UAV-RGBT
          \emph{validation} set (61,999 frames) for herding,
          since the training set was already used in Stage~1
          and may introduce data leakage. Is this valid, or
          should I use the training set and accept the
          potential leakage argument?

    \item Anti-UAV-RGBT targets are overwhelmingly small
          ($<$32\,px). If I use the three MS-COCO strata
          (small/medium/large), the small stratum will dominate
          the feature space and the herding procedure may
          select nearly identical frames. Should I use finer
          sub-strata within the small class (e.g.\ tiny $<$16\,px
          vs.\ small 16--32\,px), or use a clustering-based
          approach (e.g.\ $k$-means on embeddings within the
          small stratum) to ensure diversity?

    \item For the replay mixing ratio, I planned 10\% exemplar
          images per batch ($\approx$1--2 frames per 16-image
          batch). Is this consistent with recommendations in
          the replay literature, or should the ratio be higher
          (e.g.\ 25\%) to have a measurable effect on a
          fast-adapting model like YOLOMG?
\end{itemize}
